\section{信道容量与有噪编码定理}
\label{sec:07-Information-Theory/05-Channel-Coding/03}

噪声不会消失。接受这个前提之后，工程直觉大致给出两条路：要么认下一个正的错误率，要么不停加冗余把它压低——而上一节已经看到，重复码把错误率压向零的同时也把速率压向零。两条路都不太体面。

Shannon 给的是第三条，而且它一开始听着不像真的：\emph{存在一个正数 \(C\)，只要速率低于它，块错误概率可以做到任意小；付出的代价仅仅是速率降到 \(C\) 以下，而不是降到零。}噪声照样每次都在打击每一个符号，可整块消息可以几乎从不出错。

这条结论还有同样重要的另一半：速率一旦高过 \(C\)，错误概率就压不下去了，再精巧的编码也没用。容量因此不是一句鼓励，是一道墙。工程上知道墙在哪儿，和知道墙那边过不去，价值一样大。

\subsection{一个单字母的优化，凭什么管长块的事}

对离散无记忆信道，容量的定义短得有点可疑：

\[
C=\max_{p_X}I(X;Y).
\]

选一个单字母输入分布，算一次输入与输出之间的互信息，取最大。式子里没有块长，没有编码方案，没有错误概率。

而定理断言这个数恰好等于长块可靠通信的最高渐近速率。左边是一次概率优化，右边是一整族编码方案在 \(n\to\infty\) 时的极限行为，两者被一个等号扣在一起。信息论里最深的一步跨越就落在这条等号上，后面的论证都是在解释它为什么成立。

先把讨论对象钉住。一个 \((n,M)\) 块码把消息 \(m\in\set{1,\dots,M}\) 映成长度 \(n\) 的输入序列 \(x^n(m)\)，译码器看着收到的 \(y^n\) 猜回 \(\hat m\)，速率是 \(R=\frac1n\log_2 M\) 比特每次使用。这里的可靠不是``\(n=1000\) 时错误率低于千分之一''，而是存在一列 \(n\to\infty\) 的码，使错误概率趋于零。

\subsection{把码字当成输出空间里的球}

长块能奏效，靠的是大数定律。用信道 \(n\) 次之后，几乎所有实际出现的序列都落进``典型集''，而典型集的大小可以数出来：典型输入序列约 \(2^{nH(X)}\) 条，给定某条典型输入后与它联合典型的输出约 \(2^{nH(Y\given X)}\) 条，全部典型输出约 \(2^{nH(Y)}\) 条。这套计数在\hyperref[sec:07-Information-Theory/04-Source-Coding/04]{``AEP、典型集与信源编码定理''}里已经出现过一次，那时它给出的是压缩极限，这里给出的是通信极限。

把每个码字看成输出空间里的一个点，噪声会把它糊成一个半径由 \(H(Y\given X)\) 决定的``球''。一个球占据整个输出空间的比例大约是

\[
\frac{2^{nH(Y\given X)}}{2^{nH(Y)}}=2^{-nI(X;Y)}.
\]

于是空间里最多能塞下大约 \(2^{nI(X;Y)}\) 个互不重叠的球。码字数是 \(2^{nR}\)，什么时候塞得下、什么时候必然挤作一团，答案就写在 \(R\) 与 \(I(X;Y)\) 的大小关系里。

\begin{center}
\begin{tikzpicture}[x=1cm,y=1cm,font=\small,>=stealth]
\begin{scope}
  \draw[black!35] (0,0) rectangle (4.2,3.2);
  \foreach \c in {(0.95,0.75),(2.1,0.75),(3.25,0.75),(0.95,1.9),(2.1,1.9),(3.25,1.9)}
    {\draw[blue!45] \c circle (0.52); \fill[blue!75] \c circle (1.4pt);}
  \node[font=\footnotesize,align=center] at (2.1,-0.75) {\(R<C\)：球互不重叠\\收端几乎总能认出发的是哪个};
\end{scope}
\begin{scope}[shift={(5.9,0)}]
  \draw[black!35] (0,0) rectangle (4.2,3.2);
  \foreach \c in {(0.6,0.6),(1.35,0.75),(2.1,0.6),(2.85,0.78),(3.6,0.62),
                  (0.7,1.45),(1.5,1.55),(2.25,1.42),(3.0,1.58),(3.65,1.4),
                  (0.62,2.3),(1.4,2.4),(2.15,2.28),(2.9,2.42),(3.6,2.25),
                  (1.05,3.0),(2.5,3.0),(3.3,2.95)}
    {\draw[red!40] \c circle (0.5); \fill[red!75] \c circle (1.4pt);}
  \node[font=\footnotesize,align=center] at (2.1,-0.75) {\(R>C\)：球必然重叠\\落在交叠区的输出无法判归属};
\end{scope}
\end{tikzpicture}

\emph{方框是长度 \(n\) 的输出序列全体，点是码字，圆是它经噪声糊开后的典型输出范围。球的大小由信道定死，装不装得下只取决于你想塞多少个。容量就是这个装箱问题的临界密度。}
\end{center}

\subsection{可达性：随机地扔码字，居然就够了}

证明的骨架出人意料地不像工程。固定输入分布 \(p_X\)，随机生成约 \(2^{nR}\) 个码字，每个坐标独立按 \(p_X\) 抽——不做任何设计。译码器也简单：在码本里找唯一一个与收到的 \(y^n\) 联合典型的码字。

出错只有两种途径。真发出去的那个码字与输出不联合典型，这一项由 AEP 保证趋于零。剩下的瓶颈是冒名顶替：某个没发的码字碰巧也与输出联合典型。一个与本次传输无关的随机码字撞上这种巧合的概率在指数阶上约为 \(2^{-nI(X;Y)}\)，而竞争者有 \(2^{nR}\) 个，联合界给出

\[
2^{nR}\cdot 2^{-nI(X;Y)}=2^{-n\,(I(X;Y)-R)}.
\]

只要 \(R<I(X;Y)\)，它按指数掉到零。再把 \(p_X\) 取到最优，就得到：凡 \(R<C\)，存在一列码使错误概率趋于零。

有两处条件值得对账。其一，这是存在性论证：随机码本的\emph{平均}表现足够好，因此至少有一个具体码本不比平均差——它没有交给你任何构造方法，更没说随手生成的短码本能用。其二，平均错误概率小不等于每条消息都小，标准做法是把表现最差的一半码字扔掉，剩下的码率损失可以忽略，最大错误概率也就跟着受控。从平均到最大这一步经常被略过，但它是定理陈述的一部分。

\subsection{逆定理：墙的另一侧}

反方向的结论是：\(R>C\) 时错误概率无法趋于零。骨架同样短。假设消息均匀分布且能被高概率还原，Fano 不等式把``错误概率小''翻译成``条件熵小''，即 \(H(M\given Y^n)\le1+P_e\,nR\)；再沿 Markov 链 \(M\to X^n\to Y^n\) 用数据处理不等式，配上信道无记忆给出的 \(H(Y^n\given X^n)=\sum_i H(Y_i\given X_i)\) 与熵的次可加性 \(H(Y^n)\le\sum_i H(Y_i)\)，得到

\[
nR\le I(X^n;Y^n)+o(n)\le nC+o(n).
\]

除以 \(n\) 取极限即得 \(R\le C\)。整段论证里没有用到任何关于码的构造，所以它管住的是所有码，包括还没被发明出来的那些。

标准有限字母表信道上还有更强的版本：\(R>C\) 时最优码的错误概率趋于 1，不只是不趋于零。也就是说超过容量之后不是``稍微差一点''，而是整块消息基本必错。

\begin{center}
\begin{tikzpicture}[x=1.55cm,y=2.5cm,font=\small,>=stealth]
  \draw[->,black!55] (0,0) -- (4.5,0) node[right,font=\footnotesize] {块长 \(n\)};
  \draw[->,black!55] (0,0) -- (0,1.22) node[above,font=\footnotesize] {块错误概率};
  \draw[black!25,dashed] (0,1) -- (4.3,1);
  \node[left,font=\footnotesize] at (-0.03,1) {\(1\)};
  \node[left,font=\footnotesize] at (-0.03,0) {\(0\)};
  \draw[blue!75,thick] plot[domain=0:4.3,samples=90] (\x,{0.9*exp(-1.25*\x)});
  \draw[blue!45,thick,dashed] plot[domain=0:4.3,samples=90] (\x,{0.92*exp(-0.42*\x)});
  \draw[red!70,thick] plot[domain=0:4.3,samples=90] (\x,{1-0.62*exp(-0.85*\x)});
  \node[right,font=\footnotesize,blue!75] at (1.55,0.24) {\(R\ll C\)};
  \node[right,font=\footnotesize,blue!55] at (2.6,0.42) {\(R\lesssim C\)};
  \node[right,font=\footnotesize,red!70] at (2.5,0.86) {\(R>C\)};
\end{tikzpicture}

\emph{同一条信道，只改速率。低于容量时错误概率随块长指数下降，越靠近 \(C\) 指数越平缓；高于容量时加长块长非但无益，还会把错误概率推向 1。分界处的行为不是渐变而是翻转，这正是``硬边界''的含义。}
\end{center}

\subsection{定理承诺了什么，没承诺什么}

它给的是渐近分界线，不是性能表。块长 \(n=1000\) 时最低能做到多少错误率，临界点 \(R=C\) 上究竟能不能通信，编译码要花多少算力，时延、峰值功率、误码地板与实现精度如何影响结果——这些一概不在定理的射程内。现代编码理论和有限块长信息论做的正是补这些缺口，而结论普遍是：越想贴近容量，就得付出更长的块、更多的计算，或更精细的软信息译码。定理保证极限存在，不替工程师在有限资源下构造方案。

优化本身也有说法。互信息关于 \(p_X\) 是凹函数（\(p_Y\) 与 \(H(Y\given X)\) 都关于 \(p_X\) 线性，而熵是凹的），因此 \(\max_{p_X}I(X;Y)\) 是一个良性的凸优化问题，有限字母表下可用 Blahut--Arimoto 迭代求解。对称信道常由均匀输入达到容量，一般信道未必；不同输入的代价也可能不同，此时容量要带着预算写：

\[
C(\Gamma)=\max_{p_X:\ \E[c(X)]\le\Gamma}I(X;Y).
\]

下一节的高斯信道就是这种形式，那里的代价函数是功率。

\subsection{先压缩再纠错，为什么不吃亏}

实际系统里信源和信道是两件事：先要把数据压小，再要把压过的比特稳稳送出去。这两件事该联合设计，还是可以各干各的？

分离定理说可以各干各的。若信源每样本携带 \(H(S)\) 比特，每个样本允许使用信道 \(\kappa\) 次，则无损恢复可行的条件是

\[
H(S)<\kappa C,
\]

而且这个条件可以由``先压缩到接近熵限、再用逼近容量的码传输''这样的两段式方案达到，与任何联合设计一样好。允许平均失真 \(D\) 时把左边换成率失真函数 \(R(D)\)，结论照旧，见\hyperref[sec:07-Information-Theory/04-Source-Coding/05]{``率失真：允许失真时的压缩极限''}。

这条定理是通信系统分层的理论依据——压缩层不必知道底下跑的是光纤还是无线，传输层也不必知道自己搬的是文本还是视频。整个协议栈的模块化，很大程度上靠它撑着。

但它的前提相当具体：点对点单用户、信源与信道都无记忆、块长可以任意长、只关心无损或给定失真下的重建。这几条一旦松动，分离就可能吃亏。低时延场景下块长不能任意长，两段式方案各自的渐近优势都拿不满；多用户网络里分离一般不再最优；信道状态失配时，联合设计能提供优雅降级，而两段式方案往往在临界点上整段崩掉——这也是模拟电视信号变差时画面逐渐雪花、而数字信号常常直接黑屏的原因之一。近年把语义目标写进编码器的做法，同样是在故意放弃分离所要求的那些前提。

所以别把一条单用户渐近定理扩张成通信问题的无条件原则。它成立的地方给了工程巨大的自由，不成立的地方也很容易辨认：看块长受不受限、用户是不是只有一对、目标是不是纯粹的比特级还原。

到这里，讨论一直在离散字母表上进行。下一节把输入换成实数、把噪声换成高斯，那时``输入约束''不再是可选项——没有功率上限，容量根本没有有限值。
